\subsection{Reconstructing the margin from cash and service history}
\label{sec:insurance-retrospective-measurement}

The transition route must eventually produce a measured opening balance.
Consider a separate nonparticipating group with three equal service periods,
transition after the first, and supported zero discount rates. This example
uses the retained EU 2023/1803 provisions; it does not infer EU applicability
from the neighboring AASB or ASBE discussion. Assume no tax, foreign exchange,
OCI, incurred unpaid claims or changes in remaining expectations. The zero
rates and equal service units isolate reconstruction; they are supplied
conditions, not defaults for an insurer's portfolio
\citep[IFRS~17.C3--C8,44,B119]{eu20231803}.

The insurer starts with capital and cash 100. Before recognition it pays
acquisition costs of 10. A supported allocation assigns 6 to the recognized
group and 4 to a future group, whose acquisition asset remains recoverable.
Premium is 120, initial expected claims are 90 and initial risk adjustment
is 6. The recognized group's initial margin is consequently
\[
 120-90-6-6=18.
\]
That amount is CSM, not immediate profit. The acquisition payment appears once
in cash; the allocated 6 reduces the margin, while the future group's 4 stays
as a separate asset. These different destinations implement the acquisition
and initial-measurement requirements \citep[IFRS~17.28A--28F,38,C4]{eu20231803}.

Under full retrospective application, contemporaneous records show expected
first-period claims of 30 and risk release of 2. Actual claims of 30 are paid.
The one completed service unit releases \(18/3=6\) of CSM, leaving 12.
At transition, remaining claims 60 and risk adjustment 4 give FCF 64 and
liability \(64+12=76\). Cash is \(100-10+120-30=180\); adding the
future acquisition asset of 4 and subtracting liability gives equity 108.
Equivalently, service releases and CSM revenue, less actual claims, give
\(30+2+6-30=8\), the increase from opening capital 100.

Suppose the prior accounting basis instead reports liability 80 and acquisition
asset 10 at this same date, with the same cash 180. Prior equity is 110.
Replacing those balances gives a transition equity adjustment of
\[
 (80-76)+(4-10)=-2.
\]
There is no transition cash receipt or payment. C4(c) places the resulting net
difference in equity; it does not supply a new cash flow to make the books
balance \citep[IFRS~17.C4]{eu20231803}.

Now change the available information, keeping the baseline cash history.
Assume full retrospective application is impracticable, and that the evidence
conditions permit the specified modified-retrospective calculations. C12
reconstructs initial expected claims as remaining claims 60 plus known past
payments 30, giving 90. C14 reconstructs initial risk adjustment as the current
4 plus expected past release 2 supported by similar contracts issued at
transition. C15(b) compares one past coverage unit with two remaining units.
With the supported acquisition allocation, these inputs again yield initial
CSM 18 and transition CSM 12 \citep[IFRS~17.C6--C8,C12,C14--C15]{eu20231803}.
The agreement follows from the supplied evidence and numbers; it is not a
property of the route names.

A different actual claim payment makes the distinction visible. Suppose actual
first-period claims were 35, while the full-history initial expectation remained
90 and the remaining expectation remained 60. Full retrospection retains
CSM 12: the extra claim cost is past-service experience. Cash becomes 175 and
equity 103. In the separately admitted modified information view, C12 instead
reconstructs initial expected claims as \(60+35=95\). Initial CSM is then
\(120-95-6-6=13\), of which \(26/3\) remains. Liability is
\(64+26/3=218/3\), and equity is
\[
 175+4-\frac{218}{3}=\frac{319}{3}.
\]
The same cash supports different transition margins because the reconstruction
uses different admissible information. Replacing the full-history expectation
with the later payment would introduce hindsight; calling the modified result
the recovered historical truth would overstate what its approximation supplies.

\subsection{Prescribed missing-information treatments preserve actual cash}
\label{sec:insurance-transition-acquisition-origins}

Missing acquisition allocation has a narrower consequence than missing cash.
Return to the 30-payment case. Suppose the 10 payment is known, but reasonable
and supportable information for the C14B allocation is unavailable. C14D
prescribes nil for the specified group CSM adjustment and acquisition assets
\citep[IFRS~17.C14B--C14D]{eu20231803}. The modified initial CSM becomes
\(120-90-6=24\), leaving 16 at transition. Liability is 80 and the
acquisition asset is zero, while actual cash remains 180. Equity is therefore
100. The payment was not erased; a specified transition measurement rule changed
its asset and margin consequences. If the payment itself were unknown, this
cash reconstruction would remain unresolved.

An onerous group brings reinsurance into the same calculation. Lower the premium
to 80 and restore the supported acquisition allocation 6 and 4. Initial loss
is \(90+6+6-80=22\), with no CSM. For this unchanged, zero-rate group,
the declared systematic allocation releases one third of the loss component
with the first period's expected claim and risk releases. At transition the
loss component is \(44/3\), contained within FCF 64. It is not added to
the liability a second time. Cash is 140, so the acquisition asset of 4 gives
equity \(140+4-64=80\). C16 requires a systematic allocation under the
modified route; the full route preserves the supported historical allocation
\citep[IFRS~17.49--52,C16]{eu20231803}.

Suppose a timely treaty covers half of all these claims. Its premium is 49,
with separately supplied initial transferred risk adjustment 3 and transition
risk adjustment 2. Ordinary initial cost CSM is \(49-45-3=1\).
The initial loss recovery is \(22/2=11\), making total held cost CSM 12.
After the first of three equal service periods, held cost CSM is 8. Remaining
recovery 30 plus risk adjustment 2 gives held FCF 32 and total held asset 40.
The remaining loss-recovery component \(22/3\) sits within that asset.
This computation keeps ordinary treaty cost and recovery of the onerous loss
as separate origins \citep[IFRS~17.66--66B,B119C--B119F]{eu20231803}.

The treaty pays actual recovery 15 on the first-period claims. Combined cash
is \(140-49+15=106\), and transition equity is
\[
 106+4+40-64=86.
\]
Under the admitted modified route, C16A gives the same remaining recovery
component \((44/3)(1/2)=22/3\). Adding it to ordinary remaining cost CSM
\(2/3\) gives held cost CSM 8. Timely entry, covered-contract allocation and
the expected claims-recovery percentage remain required findings; the numerical
agreement does not establish them \citep[IFRS~17.C16A--C16B]{eu20231803}.

If the specific information needed for C16A is unavailable, C16C instead
requires that no loss-recovery component be identified. Suppose ordinary held
valuation remains supported. Held FCF is still 32 and ordinary remaining cost
CSM is \(2/3\), giving held asset \(98/3\) and equity
\(106+4+98/3-64=236/3\). This rule does not make the held asset zero or
supply a missing ordinary valuation \citep[IFRS~17.C16C]{eu20231803}.

Each calculation retains the information view, the permitted modification,
known historical cash and the separate acquisition and treaty amounts. The
reference exchange rejects hindsight in full retrospection and unsupported
approximations in modified retrospection. General discount-curve reconstruction,
direct-participation transition, changing service estimates and OCI still need
their specific procedures. These finite opening balances connect the route
conditions to measurements without treating the unresolved methods as complete.
